% Información para el motor (arara). Compilar en LuaLaTeX o XeLaTeX sólamente.

%arara: lualatex: {
%arara: --> interaction: nonstopmode,
%arara: --> options: ['-file-line-error'],
%arara: --> shell: yes,
%arara: --> synctex: yes,
%arara: --> }
%arara: bib2gls if found("aux", "glsxtr@resource")
%arara: makeglossaries if found("aux", "@istfilename")
%arara: makeindex

\documentclass[letterpaper,twoside]{article}

%%%%%%% VARIABLES ESTÁTICAS %%%%%%%
% Aqui se definen las variables para diferentes procedimientos.
% Verificar que las unidades coincidan si se deciden cambiar.

\newcommand{\defglo}[1]{\item[{\bfseries \glsname{#1}}] \glsdesc{#1}} % Define un comando para dar un item en un entorno de descripción para una palabra definida en el glosario
% \includeonly{src/PSA-2.19.tex,src/PSA-Gloss.tex}
\newcommand{\TipoID}{TEMP} % Definirlo como: PRO, PROG, ToC, G, L, CE, CI, I, etc.
\newcommand{\Prog}{TEMP}
\newcommand{\ProgTitulo}{TEMP}

\newcommand{\ie}{\emph{i.e.}}
\newcommand{\eg}{\emph{e.g.}}

%%%%%%%%%%%%%%%% Formularios %%%%%%%%%%%%%%%%
\newcommand{\IdFormAACC}{F-OP-40}

\newcommand{\BitES}{Bitácora de entradas y salidas --- Clientes (F-AD-1)}
\newcommand{\BitULl}{Bitácora de uso de llaves (F-AC-1)}
\newcommand{\BitVig}{Bitácora de vigilancia (F-AC-2)}
\newcommand{\Fpap}{Papeleta (F-AC-3)}
\newcommand{\Oent}{Orden de entrada (F-AC-4)}
\newcommand{\Osal}{Orden de salida (F-AC-5)}
\newcommand{\LayES}{Layout de carga o descarga (F-AC-6)}
\newcommand{\BitClor}{Bitácora de cloración de agua (F-AC-7)}
\newcommand{\BitTap}{Bitácora de concentración de tapete sanitario (F-AC-8)}
\newcommand{\BitUQ}{Bitácora de uso de productos químicos (F-AC-9)}
\newcommand{\BitLAd}{Bitácora de limpieza de aduana sanitaria (F-AC-10)}
\newcommand{\BitLOf}{Bitácora de limpieza de oficinas (F-AC-11)}
\newcommand{\BitLEx}{Bitácora de limpieza de exteriores (F-AC-12)}
\newcommand{\BitBPH}{Bitácora de cumplimiento de BPH (F-HyS-1)}
\newcommand{\BitCAP}{Bitácora de consumo de agua potable (F-HyS-2)}
\newcommand{\BitCE}{Bitácora de consumo de energía eléctrica (F-HyS-3)}
\newcommand{\BitIEx}{Bitácora de inspección de extintores (F-HyS-4)}
\newcommand{\BitVIR}{Bitácora de verificación de termómetros infrarrojos (F-HyS-5)}
\newcommand{\LVCMo}{Lista de verificación de condición de montacargas (F-HyS-6)}
\newcommand{\RTAC}{Registro de temperaturas con termómetro IR --- Aseguramiento de calidad (F-HyS-7)}
\newcommand{\RTMa}{Registro de temperaturas con termómetro IR --- Mantenimiento (F-HyS-8)}
\newcommand{\RAcL}{Registro de actividad de limpieza (F-HyS-9)}
\newcommand{\RAcM}{Registro de actividad de mantenimiento (F-Hys-10)}
\newcommand{\ReEq}{Registro de revisión de equipos de enfriamiento (F-MAN-1)}
\newcommand{\RCMi}{Registro de inspección de infraestructura en prevención (F-MAN-2)}
\newcommand{\Rcap}{Registro de asistencia a capacitación (F-MAN-3)}
\newcommand{\REvDQ}{Registro de eventos de derrame de químicos (F-MAN-4)}
\newcommand{\InvVyPD}{Registro de inventario de vidrio y plástico duro (F-MAN-5)}
\newcommand{\RVerSGC}{Registro de verificación del sistema de gestión de calidad (F-MAN-6)}
\newcommand{\RRCA}{Registro de recorrido de control de alérgenos (F-MC-1)}
\newcommand{\RAC}{Registro de acción correctiva o preventiva (F-OP-1)}
\newcommand{\RACQ}{Registro de acción correctiva --- Quejas (F-OP-2)}
\newcommand{\RVSCP}{Registro de verificación del servicio de control de plagas (F-OP-3)}
\newcommand{\RevD}{Registro de revisión diaria (F-SEG-1)}
\newcommand{\RCVyPD}{Registro de evento de contaminación por VyPD (F-SEG-2)}
\newcommand{\REvCA}{Registro de evento de contaminación por alérgeno (F-SEG-3)}



%%%%%%%%%%%%%%%% CÓDIGOS PSA %%%%%%%%%%%%%%%%
% \newcommand{\PSA}{PSA-3-PRO}                 % PSA-2.3 Inspección de producto durante su almacenamiento

%%%%%%%%%%%%%%%% PSA-2.3 | Inspección de producto durante su almacenamiento %%%%%%%%%%%%%%%%
\newcommand{\TiempoAndenCong}{\qty{\leq 30}{\minute}}
\newcommand{\TiempoAndenRefri}{\qty{\leq 30}{\minute}}
\newcommand{\TiempoAveriaRefri}{4}              % Tiempo máximo de permanencia de alimento refrigerado en cámara averiada (en horas) 
\newcommand{\TiempoAveriaConge}{10}             % Tiempo máximo de permanencia de alimento congelado en cámara averiada (en horas)
\newcommand{\EspacioMinimoParedProducto}{30}    % Espacio entre tarima y pared (en centimetros)
\newcommand{\EspacioMinimoTechoProducto}{30}    % Espacio entre tarima y techo (en centimetros)
\newcommand{\VecesTempManual}{dos}              % Veces que se verifica la temperatura manualmente



%%%%%%%%%%%%%%%% PSA-2.3 | Política de aprobación de proveedores %%%%%%%%%%%%%%%%
\newcommand{\NotifAuditoriaAProveedor}{2 semanas}           % Plazo máximo para notificar a proveedor sobre auditoría.
\newcommand{\PlazoResultadoAProveedor}{2 semanas}           % Plazo máximo para notificar a proveedor sobre resultados de auditoría.
\newcommand{\PlazoPlanDeAccionAuditProveedor}{2 semanas}    % Plazo maximo para recibir plan de acción de parte del proveedor tras auditoría

%%%%%%%%%%%%%%%% PSA-2.15 | Almacenamiento de registros %%%%%%%%%%%%%%%%
\newcommand{\VigenciaAlmacRegistros}{2 años}
\newcommand{\VigenciaAlmacRegistrosElec}{3 años}

%%%%%%%%%%%%%%%% GLOBAL | Resetear glosarios en cada sección %%%%%%%%%%%%%%%%
\usepackage{etoolbox}
\preto\section{\glsresetall}

\newcommand{\AC}{aseguramiento de calidad}
\newcommand{\G}{gerencia}
\newcommand{\AD}{alta dirección}
\newcommand{\OP}{operaciones}
\newcommand{\Emb}{embarques}
\newcommand{\MC}{mesa de control}	% Definición de variables estáticas
									% i.e. Tiempo máximo de cámras de refrigeración sin energía



%%%%%%%%%%%%%%%%%%%%%%%%%%% ESPECIALES %%%%%%%%%%%%%%%%%%%%%%%%%%%%%%%%%%
% \RequirePackage{expl3}    % \usepackage cannot be used before \documentclass
% \ExplSyntaxOn             % Switch on expl3 syntax
% \pdf_version_gset:n{1.7}  % Use provided expl3 function
% \ExplSyntaxOff            % Switch off expl3 syntax

%%%%%%%%%%%%%%%%%%%%%%%%%%% Tipografía %%%%%%%%%%%%%%%%%%%%%%%%%%%%%%%%%%

\usepackage[sfdefault,t,lf]{carlito}    % Alternativa OpenSource a Calibri, definir sans serif como predeterminado, numerales tabulares
% \usepackage[default,lnum]{opensans}
% \usepackage{CormorantGaramond}

% \usepackage{lucidabr}

% \usepackage{fontspec}

\fontsize{8.8pt}{12pt}\selectfont
\renewcommand{\baselinestretch}{1.2}

% \linespread{1.3}
% \DeclareUnicodeCharacter{2212}{-}

\usepackage{polyglossia}
	\setdefaultlanguage[spanishoperators=all]{spanish}  % Definir la localización del documento como es-MX (para que cuadro sea tabla, etc.)
	% \setdefaultlanguage{mexican}
	\renewcommand{\listtablename}{Índice de tablas}

\usepackage{changelog}  % Control de versiones

\usepackage{pdflscape}
\usepackage{pifont}
\newcommand{\cmark}{\ding{51}}%
\newcommand{\xmark}{\ding{55}}%

\usepackage[
	record=hybrid,% provides group and location fields (and other stuff)
	nopostdot,
	xindy,
	acronym,
	stylemods=tree]{glossaries-extra}
\setabbreviationstyle{long-short-sc}
\GlsXtrLoadResources[
	src={../glosario.bib},	% bib files
	selection={all},			% select all entries
	sort={es-MX}
]
\makeglossaries

\renewcommand{\glossarypreamble}{\fontsize{10pt}{12pt}\selectfont}
	
\usepackage{imakeidx}       % Para hacer indices
	\makeindex[columns=2]   % Estilo del indice alfabético

\newcounter{note}[section]  % Nuevo contador para notas
\renewcommand{\thenote}{\thesection.\Alph{note}}    % Definir estilo de numeración de las notas
\usepackage[framemethod=TikZ]{mdframed}             % Motor gráfico para notas (TikZ)
\newenvironment{note}[1][]{%                        % Nuevo entorno de Nota (note) 
    \refstepcounter{note}
    \begin{mdframed}[%                                  
        frametitle={Nota \thenote\ #1},                 % Titulo de las notas
        skipabove=\baselineskip{} plus 2pt minus 1pt,     % Espaciado superior
        skipbelow=\baselineskip{} plus 2pt minus 1pt,     % Espaciado inferior
        linewidth=0.5pt,                                % Tamaño de linea del recuadro
        frametitlerule=true,                            % incluir titulo
        frametitlebackgroundcolor=gray!30               % Color del recuadro
    ]%
}{%
    \end{mdframed}
}

\newcounter{contact}[section]                           % Nuevo contador para contactos
\renewcommand{\thecontact}{\thesection.\Alph{contact}}  % Definir que la numeración sea Sección.Alfabético

\newenvironment{contact}[1][]{%
    \refstepcounter{contact}
    \begin{mdframed}[%
        frametitle={Contacto \thecontact\ #1},          % Titulo del contacto
        skipabove=\baselineskip{}plus 2pt minus 1pt,	% Espaciado superior
        skipbelow=\baselineskip{}plus 2pt minus 1pt,    % Espaciado inferior
        linewidth=0.5pt,                                % Tamaño de linea del recuadro
        frametitlerule=true,                            % poner titulo
        frametitlebackgroundcolor=yellow!30             % Color del recuadro
    ]%
}{%
    \end{mdframed}
}

\usepackage[multiple]{footmisc}     % Cambios en la presentación de pies de página
\usepackage{microtype}              % Ajustes de kerning y otros 
\usepackage{graphicx}               % Mejora al motor de imagenes
\usepackage{hyperref}               % Referencias cruzadas en PDF
\usepackage{cleveref}               % Referencias cruzadas en texto
%%%%%% REDEFINICIONES DE REFERENCIAS CRUZADAS %%%%%%%%%
    \crefname{note}{nota}{notas}    
	\Crefname{note}{Nota}{Notas}
	\crefname{contact}{contacto}{contactos}
	\Crefname{contact}{Contacto}{Contactos}
	\crefname{paragraph}{parrafo}{parrafos}
	\Crefname{paragraph}{Parrafo}{Parrafos}
%%%%%%%%%%%%%%%%%%%%%%%%%%%%%%%%%%%%%%%%%%%%%%%%%%%%%%%
\newcommand{\email}[1]{\href{mailto:#1}{#1}}    % Comando para hacer que se puedan clickear emails
\newcommand{\cels}[1]{\qty{#1}{\celsius}}       % Escrir celsius rapido

\usepackage{lastpage}   % Definir variable para el numero total de páginas
\usepackage{siunitx}    % Unidades internacionales y alineamiento de decimales en tablas
	\sisetup{detect-all}    % Detectar automaticamente que la tipografía es sans serif

\usepackage{color}  % definir colores nuevos

\usepackage{tabularray} % Creación de tablas
	\UseTblrLibrary{booktabs,siunitx}   % librerias parea tablas correctas, unidades internacionales
	\definecolor{Gallery}{rgb}{0.929,0.929,0.929}   % definición del color "Gallery"
	\DefTblrTemplate{contfoot-text}{normal}{\textit{Continúa en la siguiente página}}   % Continuación de tabla
	\SetTblrTemplate{contfoot-text}{normal} % Continuación de tabla normal (footer)
	\DefTblrTemplate{conthead-text}{normal}{\textit{(Continuación)}}    % Continuación de tabla
	\SetTblrTemplate{conthead-text}{normal} % Continuación de tabla (header)

    % \usepackage[margin=1cm,bmargin=2cm,tmargin=5cm,headheight=110pt]{geometry}   % Dispoción espacial del documento
% \newgeometry{includeheadfoot, heightrounded, left=1cm, right=1cm, top=1cm, bottom=4cm, headheight=6em, footskip=1.5em}    
%     \savegeometry{L1}

% \newgeometry{includeheadfoot, heightrounded, left=1cm, right=1cm, top=1cm, bottom=4.5cm, headheight=7em, footskip=1.5em, showcrop, showframe}
%     \savegeometry{L2}

\usepackage[%
includeheadfoot,
heightrounded,
left=1cm,
right=1cm,
top=1cm,
bottom=1cm,
headheight=110pt,
footskip=1.5em]{geometry}    

\usepackage[export]{adjustbox}

\usepackage{enumitem}
	\setlist[1]{labelindent=\parindent} % Indentación como parrafo para listas
	\setlist[description]{style=nextline}
	\setlist[itemize]{leftmargin=*} 	% Eliminar margen en izq
	\setlist[itemize,1]{label=---}  	% Viñeta como em-dash
	\setlist[itemize,2]{label=--}
	\setlist{noitemsep} 				% Que no se separen los item
	% [label={\theparagraph--\arabic*}] % Numerarar parrafo-numero

	\newlist{ListaSec}{enumerate}{3}
	\setlist[ListaSec]{label* = \arabic*., name* = \arabic{section}.\arabic{subsection}}


%%%%%%%%%%%%%%%% DEFINICIÓN DE CABECERA %%%%%%%%%%%%%%%%
\usepackage{fancyhdr}   % Usar fancy headers y footers
	\renewcommand{\headrulewidth}{0pt}  % Grosor de regla del header
	\fancyhead[CE,CO,LE,LO,RE,RO]{} % Resetear headers    
    \fancyfoot[CE,CO,LE,LO,RE,RO]{} % Resetear footers

%%%%%%%%%%%%%%%%%%%%%% VARIABLES DE CADA DOCUMENTO DENTRO DE UN PROGRAMA %%%%%%%%%%%%%%%%%%%%%
	\newcommand{\Codigo}{}                  % Codigo del documento
	\newcommand{\FechaPub}{}                % Fecha de publicación
	\newcommand{\Titulo}{}                  % Titulo del documento
	\newcommand{\MayorVer}{}                % Versión mayor
	\newcommand{\MenorVer}{}                % Verisón menor
	\newcommand{\Edit}{\MayorVer.\MenorVer} % Definición de la edición

    \newcommand{\Elaboro}{}                 % Nombre de quien lo elaboró
    \newcommand{\Reviso}{}                  % Nombre de quien revisó el documento
    \newcommand{\Autorizo}{}                % Nombre de quien aprobó el documento

%%%%%%%%%%%%%%%%%% Formato por defecto para informaicón documentada general %%%%%%%%%%%%%%%%%%
\fancypagestyle{formato-PS}{%
	\fancyhead{}
	\fancyfoot{}
	\fancyhead[RE]{%
		\begin{tblr}{%	
			colspec = {X[c,m]X[c,m]},    % Especificación de columnas: Anchura estática, centrado en medio
			hlines,
			vlines,
			width = 0.35\linewidth,}%
			{\textbf{Código:}\\ \Codigo} 				 & {\textbf{Edición:}\\ \MayorVer.\MenorVer} 			\\ 	% Logotipo (2 columnas), red de frios (3 columnas), código del documento
			{\textbf{Fecha de publicación:}\\ \FechaPub} & {Página~\thepage~de~\pageref{LastPage}} 		% Titulo (4 columnas), Fecha, edición
		\end{tblr}
	}%

	\fancyhead[LO]{%
		\begin{tblr}{%	
			colspec = {X[c,m]X[c,m]},   % Especificación de columnas: Anchura estática, centrado en medio
			hlines,
			vlines,
			width = 0.35\linewidth,}%
			{\textbf{Código:}\\ \Codigo} 					& {\textbf{Edición:}\\ \MayorVer.\MenorVer} 			\\	% Logotipo (2 columnas), red de frios (3 columnas), código del documento
			{\textbf{Fecha de publicación:}\\\FechaPub}	& {Página~\thepage~de~\pageref{LastPage}} 		% Titulo (4 columnas), Fecha, edición
		\end{tblr}
	}%

	\fancyfoot[C]{Página~\thepage~de~\pageref{LastPage}}   % Para que muestre Página # de #
	\fancyfoot[L]{\textbf{CONFIDENCIAL}}                    % CONFIDENCIAL
	\fancyfoot[R]{\Codigo~v.\Edit}                          % Estructura del código CODIGO v.EDICION
}%

\fancypagestyle{formato-PI}{%
	\fancyhead{}
	\fancyfoot{}
    \fancyhead[C]{%
	    \begin{tblr}{%	
            colspec = {X[c,m]X[c,m]X[c,m]X[c,m]X[c,m]X[c,m]},    % Especificación de columnas: Anchura estática, centrado en medio
            hlines,
            vlines,
			width = \linewidth,}%
			\SetCell[c=2]{} \includegraphics[height=3em, valign=c]{./RDF_Logo.png}	&	& \SetCell[c=3]{} \large{Red de fríos S.A. de C.V.}	&	&												& {\textbf{Código:}\\ \Codigo}	\\ % Logotipo (2 columnas), red de frios (3 columnas), código del documento
			\SetCell[c=4]{} {\textbf{Titulo:}\\ \Titulo} 							&	&													&	& {\textbf{Fecha de publicación:}\\ \FechaPub}	& {\textbf{Edición:}\\ \Edit}	\\ % Titulo (4 columnas), Fecha, edición
            \end{tblr}
	}
    \fancyfoot[C]{Página~\thepage~de~\pageref{LastPage}}   % Para que muestre Página # de #
    \fancyfoot[L]{\textbf{CONFIDENCIAL}}                    % CONFIDENCIAL
    \fancyfoot[R]{\Codigo~v.\Edit}                          % Estructura del código CODIGO v.EDICION
}%

%%%%%%%%%%%%%%%%%%%%%%%%%%%%%%%%%%%%%%%%%%%%%%%%%%%%%%%%%%%%%%%%%%%%%%%%%%%%%%%%%%%%%%%%%%%%%%%%%%%%%%%%%%%%%%%%%%
			
\setcounter{tocdepth}{4}            % Qué tanto mostrará la tabla de contenidos
\setcounter{secnumdepth}{6}         % Que tantas secciones numerará LaTeX (Se ajustó hasta subparrafo)

 % Define un comando para dar un item en un entorno de descripción para una palabra definida en el glosario

% \includeonly{src/BPM-1.1.tex}
\renewcommand{\Prog}{VyPQ}

\begin{document}
\pagestyle{formato-PS}
\thispagestyle{formato-PI}
\renewcommand{\MayorVer}{1}
\renewcommand{\MenorVer}{0}
\renewcommand{\TipoID}{ToC}
\renewcommand{\Codigo}{\Prog--\TipoID}
\renewcommand{\FechaPub}{2023--01}
%\renewcommand{\Edit}{01}
\renewcommand{\Titulo}{Programa de control de vidrio y plástico quebradizo --- Índice}

\tableofcontents
\thispagestyle{formato-PI}
\renewcommand{\MayorVer}{2}
\renewcommand{\MenorVer}{0}
\renewcommand{\Titulo}{Programa de control de vidrio y plástico quebradizo}
\renewcommand{\TipoID}{PRO}
\renewcommand{\FechaPub}{2023--01}
% \renewcommand{\Titulo}{Politica de calidad}

\section{\Titulo}\index{Programa!control de!vidrio y plástico quebradizo}
\renewcommand{\Codigo}{\Prog--\thesection--\TipoID}

\subsection{Objetivo}
Establecer procedimiento de verificación y mantenimiento de los artículos o instalaciones que contengan vidrio o plástico quebradizo para evitar o minimizar un posible riesgo de contaminación del producto por estos artículos.

\subsection{Alcance}
Todo material de vidrio o plástico quebradizo ubicados dentro del almacén.

\subsection{Terminología y definiciones}
\begin{description}
	\defglo{limpieza-y-sant}
\end{description}

\subsection{Documentos y/o normas relacionadas}
\begin{itemize}
	\item Programa de Buenas Prácticas de Distribución;
	\item Programa de Seguridad Alimentaria.
\end{itemize}

\subsection{Procedimiento}
El desarrollo contempla que el vidrio requiere estar protegido dentro de la planta, cuáles son las actividades que se tienen que realizar si se presenta una ruptura de vidrio o plástico quebradizo dentro de las áreas de Operación y como se direcciona el almacenamiento y uso de empaque de vidrio o plástico quebradizo. Además de cómo prevenir la contaminación por ruptura o astillamiento.
El programa contempla:

\begin{itemize}
	\item Frecuencia de verificación de control de vidrio y plástico quebradizo o quebradizo.
	\item Limpieza de áreas y equipo por contaminación de vidrio y plástico quebradizo.
	\item Procedimiento de que hacer en caso de contaminación de equipo, área, producto y material de empaque.
	\item Listado de vidrio y plástico quebradizo por área y mapas de su ubicación
	\item Lista de verificación.
	\item Todos los vidrios mencionados en la lista, están protegidos para evitar astillamiento o ruptura y así evitar contaminación del producto.
	\item Algunos ya son de fábrica inastillable, como todos los vidrios de las ventanas y lámparas del área de Almacenamiento.
	\item Las que no se encuentran en esta categoría deberán estar recubiertas con plástico adherente, capuchones plásticos o protectores (cubre polvos) para evitar los astillamientos en caso de ruptura.
	\item La verificación del control de vidrio y plástico quebradizo en las áreas de operación se realiza de manera MENSUAL por parte de parte del personal de mantenimiento.
	\item En caso de cualquier evento este se registrará en el momento ocurrido.
	\item El supervisor de almacén deberá informar inmediatamente de vidrio roto o plástico que sea detectado en el almacén.
	\item El supervisor de almacén debe inspeccionar la rotura para determinar si algún producto o empaque de alimentos ha sido contaminado. Con base en las observaciones del Supervisor de Almacén, se determinará si el producto puede ser salvado o, si está contaminado, debe ser enviado a cuarentena y se la dará aviso al cliente del evento sucedido
	\item En caso de ruptura se procederá al procedimiento de limpieza de áreas y equipo por contaminación de Vidrio y plástico quebradizo.
	\item\textbf{NOTA:} Se anexa el formato de inventario por áreas con el contenido de materiales de vidrio y plástico quebradizo, así como el mapa de la planta, con la ubicación en cada una de las áreas.
\end{itemize}

\subsection{Acciones correctivas}
\begin{itemize}
	\item Cuando se presente una no conformidad en la realización del procedimiento detectada por el supervisor o encargado del área se deberá repetir este mismo.
	\item Registrar acciones correctivas correspondientes.
\end{itemize}

\subsection{Frecuencia}
\begin{itemize}
	\item Mensualmente se llena la lista de verificación de materiales de vidrio y plástico quebradizo;
	\item cada evento por contaminación de vidrio o plástico quebradizo.
\end{itemize}


\begin{changelog}[title=Registro de cambios,simple, sectioncmd=\subsection*,label=changelog-\thesection-\MayorVer.\MenorVer]
	\begin{version}[v=\MayorVer.\MenorVer, date=2023--01, author=Pablo E. Alanis]
		\item Cambio de formato;
		\item Cambios en la serialización de versiones;
		\item Correcciones ortográficas y de estilo.
	\end{version}

	\begin{version}[v=1.7, date=2022--05, author=Alonso M.]
		\item cambio de fecha;
	\end{version}

	\shortversion{v=1.6, date=2021--05, changes=No hubo cambios}
\end{changelog}


\thispagestyle{formato-PI}
\renewcommand{\MayorVer}{2}
\renewcommand{\MenorVer}{0}
\renewcommand{\Titulo}{Lavado de áreas y equipos por contaminación de vidrio y plástico quebradizo}
\renewcommand{\TipoID}{IT}
\renewcommand{\FechaPub}{2023--01}

\section{\Titulo}\index{Programa!control de!vidrio y plástico quebradizo!limpieza en caso de contaminación por VyPQ}
\renewcommand{\Codigo}{\Prog--\thesection--\TipoID}

\subsection{Objetivo}
Establecer procedimiento que asegure un mantenimiento y limpieza apropiados, manejo de los desechos y vigilancia de la efectividad de los procedimientos de las áreas, equipos y producto que estén involucrados en contaminación con vidrio y/o plástico quebradizo.

\subsection{Alcance}
Todas las áreas afectadas por partículas de vidrio y plástico quebradizo.

\subsection{Terminología y definiciones}
\begin{description}
	\defglo{limpieza-y-sant}
\end{description}

\subsection{Procedimiento}

\subsubsection{Materiales}

\begin{itemize}
	\item Garrafa identificada con agua y detergente a la concentración especificada por el proveedor.
	\item Cepillos para equipo de color rojo con franjas naranja.
	\item Cepillo para piso (escoba) color rojo con franjas naranja.
	\item Recogedor rojo con franjas naranja.
	\item Aspiradora.
	\item Bolsa de basura.
	\item Papel absorbente desechable.
	\item Contenedor para desperdicio.
\end{itemize}

\subsubsection{Instrucciones}

\paragraph{Contaminación de equipo y área}

\begin{itemize}
	\item Limpie el equipo.
	\item Lave el equipo (SEGÚN EL PROCEDIMIENTO DE LAVADO DE CADA EQUIPO)
	\item Lave o limpie los equipos que se encuentran en un radio de 15 metros (CADA EQUIPO SEGÚN EL PROCEDIMIENTO DE LIMPIEZA DE CADA UNO DE ELLOS)
	\item Lavar estanteros y todo inmobiliario que se encuentre en un radio de 5 metros.
	\item Barra con escoba en el piso en un radio de 5 metros.
	\item Junte basura con recogedor color rojo con franja naranja
	\item Coloque en bolsa para basura. (eliminar al momento y cambie la bolsa de basura)
	\item Aplicar agua con jabón en todas las partes a lavar del piso
	\item Cepille muy bien el piso con cepillo de piso color rojo con franja naranja.
	\item Enjuague con agua corriente.
\end{itemize}

\paragraph{De vidrio y/o plástico quebradizo}

\begin{itemize}
	\item En caso de evaluar un proceso de contaminación MUY GRAVE, separa el producto al área de rechazo e identifique el producto con el problema.
	\item Se tendrá comunicación con el cliente.
	\item Se procederá según las indicaciones del departamento de calidad del cliente.
	\item (Elimine el producto dañado)
	\item A continuación se procede a la limpieza de equipos y áreas
	\item Elimine material de empaque que se localice en un radio de 5mts (este se aparta y se manda al área de materiales no conformes)
	\item Lave el área (SEGUN EL PROCEDIMIENTO DE LIMPIEZA DE CADA EQUIPO)
	\item Lave o limpie los equipos que se encuentran en un radio de 5 metros (CADA EQUIPO SEGÚN EL PROCEDIMIENTO DE LIMPIEZA DE CADA UNO DE ELLOS)
	\item Lavar estanteros y todo inmobiliarios que se encuentre en un radio de 5 metros.
	\item Barra con escoba el piso en un radio de 5 metros.
	\item Junte basura con recogedor rojo con franja naranja.
	\item Coloque en bolsa para basura. (Eliminar al momento y cambie la bolsa de basura)
	\item Aplicar agua con jabón en todas las partes a lavar del piso
	\item Cepille muy bien el piso con cepillo de piso color rojo con franja naranja.
	\item Enjuague con agua corriente.
	\item Elimine el exceso de agua con jalador para piso color rojo con franja naranja.
	\item Recoger todos los artículos que fueron utilizados para ser acomodados en el área que corresponde.
	\item Avise al jefe inmediato que el trabajo ya está para su supervisión
	\item Registre el evento en el formato “registro de eventos de contaminación con vidrio y plástico quebradizo” numero de orden de mantenimiento si requiere cambio o reparación y el procedimiento de limpieza realizado de acuerdo al programa.
\end{itemize}

\subsection{Responsables de la actividad}

NOTA: cada uno se analizara y se tomara las acciones correctivas según su gravedad. Las personas autorizadas para valorar esta gravedad (Si la presencia lo permite la decisión se toma en conjunto) son:

\begin{itemize}
	\item Gerente de Operaciones.
	\item Cliente
	\item Jefe de Almacén
	\item Gerente General
\end{itemize}

\subsection{Acciones correctivas}

\begin{itemize}
	\item Cuando se presente una no conformidad en la realización del procedimiento, detectada por el supervisor o encargado del área se deberá de repetir este mismo hasta que sea corregida la desviación.
	\item Registrar en formato de acciones correctivas correspondientes.
\end{itemize}

\subsection{Frecuencia}

\begin{itemize}
	\item Cuando se presente evento de rotura de vidrio y/o plástico quebradizo.
\end{itemize}

\begin{changelog}[simple, sectioncmd=\subsection*,label=changelog-\thesection-\MayorVer.\MenorVer]
	\begin{version}[v=\MayorVer.\MenorVer, date=2023--01, author=Pablo E. Alanis]
		\item Cambio de formato;
		\item Cambios en la serialización de versiones;
		\item Correcciones ortográficas y de estilo.
	\end{version}

	\begin{version}[v=1.7, date=2022--05, author=Alonso M.]
		\item cambio de fecha;
	\end{version}

	\shortversion{v=1.6, date=2021--05, changes=No hubo cambios}
\end{changelog}


\renewcommand{\MayorVer}{1}
\renewcommand{\MenorVer}{0}
\renewcommand{\TipoID}{G}
\renewcommand{\Codigo}{\Prog-\TipoID}
\renewcommand{\FechaPub}{2023--01}
%\renewcommand{\Edit}{01}
\renewcommand{\Titulo}{Programa de control de Vidrio y Plástico Quebradizo --- Glosario}

\printunsrtglossary[type=acronym,nonumberlist]
\printunsrtglossary[type=\glsdefaulttype,nonumberlist]
% \printunsrtglossary[type=main]    
% \printglossary[type=acronym]
% \printglossary[type=main]
% \printunsrtglossary[type=acronym]
% \printunsrtglossary[type=main]
\renewcommand{\MayorVer}{1}
\renewcommand{\MenorVer}{0}
\renewcommand{\TipoID}{I}
\renewcommand{\Codigo}{\Prog-\TipoID}
\renewcommand{\FechaPub}{2023--01}
%\renewcommand{\Edit}{01}
\renewcommand{\Titulo}{Indice alfabético}
\printindex
\end{document}